% Section 5 only. Reference: blackboard_feedback_report.tex in
% docs/theory/blackboard_theory/blackboard_feedback_report_with_available_susceptibility_source.zip
\section{Efficiency of Feedback Control}
\label{sec:feedback_efficiency}

We use a solvable blackboard reference to separate sensing, message exposure,
and behavioral response. Let $n_k\in\{0,\ldots,N\}$ count supporters of the
controller target $Z$ before round $k$, and write $x_k=n_k/N$. One feedback
cycle has causal order $n_k\to Y_k\to U_k\to n_{k+1}$, where $U_k=0$ denotes
silence and $U_k=1$ active posting. Both branches contain $M$ population
updates, usually $M=N$. Posting $b$ messages changes the chance of exposure;
it does not guarantee $b$ controlled revisions. The results below are exact
for the reference kernels, rather than a fitted description or an energetic
thermodynamics of the full language-model system.

\subsection{Sensing, posting, and sampled exposure}

The controller senses $q_c$ agents without replacement. For $Y=y$ sampled
target supporters, the sensor and soft feedback policy are
\begin{align}
S(y\mid n)&=\frac{\binom{n}{y}\binom{N-n}{q_c-y}}{\binom{N}{q_c}},
\label{eq:sensor_kernel_single_affinity}\\
\pi(1\mid y)&=\sigma\!\left[\beta\left(\theta-\frac{y}{q_c}\right)\right],
\qquad \sigma(z)=(1+e^{-z})^{-1},
\label{eq:feedback_policy_single_affinity}\\
a_n&=P(U=1\mid n)=\sum_y S(y\mid n)\pi(1\mid y).
\label{eq:advocacy_probability_single_affinity}
\end{align}
Suppose the frozen eligible board contains $B$ peer messages and, when active,
$b$ controller messages. Each focal update samples $s=\min(q,B+b)$ messages
uniformly without replacement. Its probability of seeing at least one
controller message is
\begin{equation}
\boxed{w_b=1-\frac{\binom{B}{s}}{\binom{B+b}{s}}.}
\label{eq:blackboard_exposure}
\end{equation}
Here $\binom{B}{s}=0$ when $s>B$. The theory uses the number of messages
\emph{actually posted}; a nominal budget must be replaced by the realized
count when fewer messages are posted. Thus $q_c,\theta,\beta$ govern the
posting decision, while $B,b,q$ govern exposure. For positive $q$, exposure
saturates as $b$ grows.

\subsection{From one agent update to a full feedback round}

\paragraph{One update: willingness to revise and preference after revision.}
At each microscopic update, one agent is selected uniformly from the $N$
agents. We describe its response using two probabilities. The parameter
$\gamma\in[0,1]$, called \emph{kinetic compliance}, is the probability that
the agent revises its vote. Conditional on revising, $p\in[0,1]$ is the
probability that it selects the controller target $Z$; with probability
$1-p$ it selects a non-target alternative. With probability $1-\gamma$,
it keeps its current vote. Revision here means making a new choice, which
can agree with the old one. Thus $\gamma$ is not the probability of an
observed vote change. For example, $\gamma=0.4$ and $p=0.8$ mean a
$40\%$ chance of making a new choice and, conditional on that choice, an
$80\%$ chance of selecting $Z$.

\paragraph{A kernel is a table of transition probabilities.}
Write $K_{\gamma,p}(m\mid n)$ for the probability that this single update
leaves $m$ target supporters, given $n$ initially. The symbol $K$ denotes
this transition table, or \emph{kernel}; the subscripts specify which
revision probabilities it uses. Only one agent updates, so the count can
increase by one, decrease by one, or stay fixed:
\begin{align}
K_{\gamma,p}(n+1\mid n)&=\gamma(1-n/N)p,\\
K_{\gamma,p}(n-1\mid n)&=\gamma(n/N)(1-p),\\
K_{\gamma,p}(n\mid n)&=1-K_{\gamma,p}(n+1\mid n)-K_{\gamma,p}(n-1\mid n).
\label{eq:single_affinity_kernel}
\end{align}
For the first line, the selected agent must be a non-supporter, an event
of probability $1-n/N$, then revise with probability $\gamma$, and choose
$Z$ with probability $p$. The second line reverses this reasoning: a
supporter is selected, revises, and chooses a non-target alternative.
The last line collects all ways to leave the supporter count unchanged.

\paragraph{Two response classes: without and with exposure.}
An agent may respond differently when its sampled messages include
controller content. We therefore use the same transition rule with two
parameter pairs:
\[
K_0=K_{\gamma_0,p_0},\qquad K_1=K_{\gamma_1,p_1}.
\]
The subscript $0$ means \emph{no controller message was sampled}:
$\gamma_0$ is the baseline revision probability and $p_0$ is the baseline
probability of choosing $Z$ after revision. The subscript $1$ means
\emph{at least one controller message was sampled}: $\gamma_1$ and $p_1$
are the corresponding exposed-response probabilities. These are effective
behavioral parameters, not message counts. Exposure may change willingness
to revise, preference for the target, or both. In particular, $p_0$ need
not be zero: peer messages can support $Z$ without controller exposure.

\paragraph{An active controller does not expose every update.}
When the controller posts $b$ messages, a focal update sees controller
content with probability $w_b$ from Eq.~\eqref{eq:blackboard_exposure}.
It therefore follows $K_1$ with probability $w_b$ and $K_0$ otherwise.
The resulting single-update kernel is their probability-weighted mixture,
\begin{equation}
K_b=(1-w_b)K_0+w_bK_1=K_{\gamma_b,p_b}.
\end{equation}
Here the subscript $b$ refers to the number of posted messages, whereas
$0$ and $1$ above label exposure classes. The mixture has the same form as
the original rule, with effective parameters
\begin{align}
\gamma_b&=(1-w_b)\gamma_0+w_b\gamma_1,\\
p_b&=\frac{(1-w_b)\gamma_0p_0+w_b\gamma_1p_1}{\gamma_b}.
\label{eq:blackboard_mixture}
\end{align}
The first expression is the overall probability of revision. The numerator
of the second is the probability of revising and choosing $Z$; dividing by
$\gamma_b$ gives the preference conditional on revision. Thus $p_b$ weights
each exposure class by how often it produces revision, not just how often
it is encountered. If $\gamma_b=0$, no revision occurs, $K_b$ leaves every
state unchanged, and $p_b$ is unnecessary.

\paragraph{From one update to a round.}
The kernels $K$ describe one agent update. We use $Q$ for the transition
table of a full round containing $M$ successive updates. Repeating the
same update rule $M$ times gives a matrix power:
\begin{equation}
\boxed{Q_0=K_0^M,\qquad Q_1=K_b^M.}
\label{eq:controlled_layer_kernels}
\end{equation}
For $Q$, the subscript labels the \emph{controller action}: $Q_0$ is the
silent-round kernel and $Q_1$ the active-round kernel. An active round uses
$K_b$, not $K_1$, because individual updates may still miss the controller
messages. A silent round uses $K_0$ throughout, but can still change votes
through baseline social interaction. Finally, before the controller action
is known, the round kernel averages over its posting probability $a_n$:
\begin{equation}
Q_\pi(m\mid n)=(1-a_n)Q_0(m\mid n)+a_nQ_1(m\mid n).
\end{equation}
Thus the notation follows three steps: $K_0,K_1$ specify responses to
exposure; $K_b$ averages over whether exposure occurs; and $Q_0,Q_1,Q_\pi$
describe full rounds, first conditional on the controller action and then
averaged over the policy. The matrix-power construction assumes a frozen
board, common response parameters across agents and updates, and the same
unexposed baseline in both controller branches.

\subsection{Raw and availability-normalized susceptibility}
\label{sec:control-susceptibility}

At a common initial state, the primary response observable is
\begin{equation}
\boxed{\chi_b(x)=\mathbb{E}[x'\mid U=1,x]-\mathbb{E}[x'\mid U=0,x].}
\label{eq:control-susceptibility}
\end{equation}
The one-update mean obeys
$\mathbb{E}[x_{j+1}\mid x_j]=(1-\gamma/N)x_j+\gamma p/N$.
Consequently the exact round mean and susceptibility are
\begin{align}
F_M(x;\gamma,p)&=p+(x-p)(1-\gamma/N)^M,\\
\chi_b(x)&=F_M(x;\gamma_b,p_b)-F_M(x;\gamma_0,p_0)\nonumber\\
&=p_b(1-r_b)-p_0(1-r_0)+(r_b-r_0)x,
\label{eq:chi_finite_bias}
\end{align}
where $r_b=(1-\gamma_b/N)^M$ and $r_0=(1-\gamma_0/N)^M$.
The reference predicts an affine dependence on initial target share, with
budget entering through exposure. If $r_b\ne r_0$, the zero-response point is
$x_b^*=[p_0(1-r_0)-p_b(1-r_b)]/(r_b-r_0)$; it need not lie in $[0,1]$.
When $\gamma_0=\gamma_1=\gamma$, the state dependence cancels:
\begin{equation}
\chi_b(x)=w_b(p_1-p_0)\left[1-(1-\gamma/N)^M\right].
\label{eq:blackboard_equal_compliance}
\end{equation}
Exposure and compliance therefore impose distinct limits: being seen more
often and changing behavior after exposure are separate mechanisms.

To express the same branch contrast relative to the remaining non-target
population, define
\begin{equation}
\boxed{\chi_{b,\mathrm{avail}}(x)=\frac{\chi_b(x)}{1-x},\qquad x<1.}
\label{eq:blackboard_available_susceptibility}
\end{equation}
The saturated state $x=1$ is excluded. This diagnostic can be negative or
exceed one: the contrast includes differential retention of existing target
supporters as well as additional adoption. It is neither a conversion
probability nor a bounded efficiency. For occupancy $p(n)$, the corresponding
available-mass-weighted summary is
\begin{equation}
\overline{\chi}_{\mathrm{avail}}[p]
=\frac{\sum_{n=0}^{N-1}p(n)\chi_b(n/N)}
{\sum_{n=0}^{N-1}p(n)(1-n/N)},
\label{eq:blackboard_available_summary}
\end{equation}
when the denominator is positive. Both sums exclude saturation; this ratio
avoids an unweighted average of potentially large near-saturation ratios.

\subsection{Action information and information--response efficiency}

The exact state-local action-to-population information is the weighted
Jensen--Shannon divergence of the two evolving round kernels:
\begin{equation}
T_\pi(n)=I(U;n'\mid n)
=\sum_u P(u\mid n)\sum_m Q_u(m\mid n)
\log_2\frac{Q_u(m\mid n)}{Q_\pi(m\mid n)}.
\label{eq:action-population-information}
\end{equation}
With $H_2(a)=-a\log_2a-(1-a)\log_2(1-a)$,
$0\le T_\pi(n)\le H_2(a_n)$. Thus the action-information fraction is
\begin{equation}
\eta_{\mathrm{IF}}(n)=\frac{T_\pi(n)}{H_2(a_n)}\in[0,1],
\label{eq:blackboard_eta_if}
\end{equation}
when $H_2(a_n)>0$. It measures how much conditional action uncertainty is
visible in the next state, rather than directional motion.

Because $x'\in[0,1]$, Pinsker's inequality gives
\begin{equation}
T_\pi(n)\ge B_{\mathrm{IR}}(n)
\equiv\frac{2a_n(1-a_n)}{\ln2}\chi_b(n/N)^2.
\label{eq:information-response-bound}
\end{equation}
The bounded information--response efficiency is therefore
\begin{equation}
\boxed{\eta_{\mathrm{IR}}(n)
=\frac{2a_n(1-a_n)\chi_b(n/N)^2}{(\ln2)T_\pi(n)}\le1,}
\label{eq:eta_ir}
\end{equation}
for $T_\pi(n)>0$. This bound uses the \emph{raw} susceptibility, not its
availability-normalized counterpart, and does not require $Q_0$ to be the
identity. An informative action channel can have little mean response if
its effect is mainly on fluctuations or higher moments.

\subsection{Excess current and finite-time path accounting}
\label{sec:thermodynamic_efficiency}

Let $p_k(n)$ denote the pre-round ensemble and
$\mu_u(n)=\mathbb{E}[n'-n\mid n,U=u]$. Controller-excess current subtracts the
silent counterfactual:
\begin{equation}
\boxed{J_{c,k}^{\mathrm{ex}}
=\sum_n p_k(n)a_n[\mu_1(n)-\mu_0(n)]
=N\sum_n p_k(n)a_n\chi_b(n/N).}
\label{eq:mean_controlled_current_single_affinity}
\end{equation}
Equivalently, $J_{c,k}^{\mathrm{ex}}=J_{\mathrm{tot},k}-J_{\mathrm{sil},k}$,
where $J_{\mathrm{tot},k}=\sum_np_k(n)[(1-a_n)\mu_0(n)+a_n\mu_1(n)]$
and $J_{\mathrm{sil},k}=\sum_np_k(n)\mu_0(n)$.

For interior preferences, define $h_0=\ln[p_0/(1-p_0)]$ and
$h_b=\ln[p_b/(1-p_b)]$, with $h_u=h_0$ under silence and $h_u=h_b$
under active posting. On supported transitions, weighted reversibility gives
\begin{equation}
\ln\frac{Q_u(m\mid n)}{Q_u(n\mid m)}
=\ln\binom{N}{m}-\ln\binom{N}{n}+h_u(m-n).
\label{eq:blackboard_branch_ratio}
\end{equation}
Let $p_{Y,k}(y)=\sum_np_k(n)S(y\mid n)$ and let $p_{k+1}$ be the ensemble
propagated by $Q_\pi$. The forward and reverse-reference path measures are
\begin{align}
P_F(n,y,u,m)&=p_k(n)S(y\mid n)\pi(u\mid y)Q_u(m\mid n),\\
P_R(m,y,u,n)&=p_{k+1}(m)p_{Y,k}(y)\pi(u\mid y)Q_u(n\mid m).
\end{align}
Their log ratio yields the exact finite-time identity
\begin{equation}
\boxed{\Sigma_k=D_{\mathrm{KL}}(P_F\Vert P_R)
=\Delta S_{\mathrm{sys},k}+I(n_k;Y_k)+\mathcal{A}_k\ge0,}
\label{eq:finite_time_second_law}
\end{equation}
where all terms use natural logarithms and
\begin{align}
S_{\mathrm{sys}}[p]&=-\sum_n p(n)\ln p(n)
+\sum_n p(n)\ln\binom{N}{n},
\label{eq:system_entropy_functional}\\
\mathcal{A}_k&=\sum_n p_k(n)
\left[(1-a_n)h_0\mu_0(n)+a_n h_b\mu_1(n)\right].
\label{eq:blackboard_affinity_flow}
\end{align}
The entropy change is retained; stationarity is not assumed.

For $\mathcal{A}_k\ge0$ and positive denominator, the bounded
\emph{total path-transport fraction} is
\begin{equation}
\eta_{\mathrm{path},k}
=\frac{\mathcal{A}_k}{\Sigma_k-\Delta S_{\mathrm{sys},k}}
=\frac{\mathcal{A}_k}{\mathcal{A}_k+I(n_k;Y_k)}\in[0,1].
\label{eq:blackboard_eta_path}
\end{equation}
It includes autonomous baseline transport and is not a controller-only
thermodynamic efficiency. Indeed, with
$J_{\mathrm{adv},k}=\sum_np_k(n)a_n\mu_1(n)$,
\begin{equation}
\mathcal{A}_k-h_0J_{\mathrm{sil},k}
=h_0J_{c,k}^{\mathrm{ex}}+(h_b-h_0)J_{\mathrm{adv},k}.
\end{equation}
This excess-affinity diagnostic need not be nonnegative and is not itself a
second-law-bounded expenditure. A sufficient attribution condition is a
neutral baseline affinity, $p_0=1/2$ or $h_0=0$. Only under this additional
condition does the path fraction become
\begin{equation}
\eta_{\mathrm{th},k}^{(h_0=0)}
=\frac{h_bJ_{\mathrm{adv},k}}
{h_bJ_{\mathrm{adv},k}+I(n_k;Y_k)}\in[0,1],
\label{eq:eta_th}
\end{equation}
provided $h_bJ_{\mathrm{adv},k}\ge0$ and the denominator is positive.
Neutral affinity does not imply zero silent drift. In general, the pair
$(J_{c,k}^{\mathrm{ex}},\Sigma_k)$ describes excess steering versus total coarse
path irreversibility without assigning baseline transport to the controller.

\paragraph{Empirical scope.}
Exposure frequency and exposed/unexposed transition rates can test the
reference. If $p_{+,r}$ and $p_{-,r}$ are conditional target-entry and exit
probabilities in response class $r$, then
$\gamma_r=p_{+,r}+p_{-,r}$ and
$p_r=p_{+,r}/(p_{+,r}+p_{-,r})$ when the denominator is nonzero.
Shared parameters should predict held-out budgets and states; observational
exposure classes alone do not establish a causal effect. Paired active and
silent continuations should match the complete state, including board contents
and epistemic information. Board growth, self-exclusion, message content, and
heterogeneous epistemic state can violate the frozen, homogeneous reference.
Such departures call for explicit conditioning or time-dependent kernels.
The safe bounded information quantities are $\eta_{\mathrm{IF}}$ and
$\eta_{\mathrm{IR}}$; controller-specific thermodynamic attribution requires
additional assumptions.
